% !PRINTED 566
\weberpage{566}

\begin{center}
§ 151. \textbf{Die komplexe Multiplikation der Funktion} $\wp(u)$.
\end{center}

Wir nehmen jetzt an, die Perioden $\omega_1$, $\omega_2$ der Funktion $\wp(u)$ genügen einer quadratischen Gleichung
\weq{1}{A\omega_2^2 + B\omega_2\omega_1 + C\omega_1^2 = 0,}
in der $A$, $B$, $C$ ganze rationale Zahlen ohne gemeinschaftlichen Teiler sind. Es sei $A$ positiv, und die Diskriminante der quadratischen Form $(A, B, C)$
\weq{2}{B^2 - 4AC = \Delta}
sei eine \gesperrt{negative Stammdiskriminante}, d.\,h. es enthalte $\Delta$ keinen quadratischen Faktor, nach dessen Absonderung eine Diskriminante übrig bleibt (§ 84).

Setzen wir
\weq{3}{\omega = \dfrac{\omega_2}{\omega_1},}
so folgt aus (1):
\weq{4}{2A\omega = -B + \sqrt{\Delta},}
und wenn wir $\sqrt{\Delta}$ positiv imaginär annehmen, so erhält $\omega$ einen positiv imaginären Teil und kann als Modul einer $\vartheta$-Funktion dienen. $\Delta$ ist die Grundzahl eines imaginären quadratischen Körpers, den wir mit $\Omega$ bezeichnen\textsuperscript{1)}\begingroup\renewcommand{\thefootnote}{1)}\footnotetext[1]{Die Annahme, daß $\Delta$ Stammdiskriminante sei, ist hier zur Vereinfachung gemacht. In meiner Abhandlung \glqq Über Zahlgruppen in algebraischen Körpern\grqq{} (Mathematische Annalen, Bd. L) ist diese Annahme nicht gemacht. Man erhält dann allgemeinere Körper, die zu den Ordnungen gehören, wie in § 124.}\endgroup.

Durch die singuläre Invariante $j(\omega)$ wird der Klassenkörper $\mathfrak{R}(\Delta)$ bestimmt, dessen Relativgrad in bezug auf $\Omega$ gleich der Klassenzahl der Diskriminante $\Delta$ ist. Die ganzen Zahlen des quadratischen Körpers $\Omega = \mathfrak{R}(\sqrt{\Delta})$ sind von der Form:
\weq{5}{\mu = \dfrac{x + y\sqrt{\Delta}}{2},}
worin $x$ und $y$ ganze rationale Zahlen sind, die der Bedingung
\[ x \equiv By \pmod{2} \]
genügen. Zu der durch (5) definierten Zahl $\mu$ bestimmen wir vier ganze rationale Zahlen $a$, $b$, $c$, $\partial$:
\weq{6}{\begin{aligned} a &= \dfrac{x+By}{2}, & b &= Ay, \\ c &= -Cy, & \partial &= \dfrac{x-By}{2}, \end{aligned}}
%CONT>>
